\section{选题目标的重新定义}

在前两部分分析之后，可以看到，当前问题的关键不在于系统是否已经足够复杂，也不在于是否还能继续增加功能，而在于现有选题的研究对象需要重新定义。如果这一点不改变，那么后续所有改进都很容易继续停留在“把生命科学多智能体助手做得更完善”这一层，而难以形成一个边界清晰、创新点集中、能够独立成立的毕业设计课题。

因此，下一步最重要的不是继续扩展原系统，而是完成从“应用系统叙事”到“研究问题叙事”的转换。

从原有思路看，课题更容易被理解为：构建一个面向生命科学任务，尤其是蛋白设计任务的多智能体系统，通过规划、执行、工具调用和人工介入，帮助完成相关科研流程。这种表述的主语仍然是“生命科学助手”。一旦如此，无论其内部加入多少运行机制，都仍然可能被视为更大助手系统中的一个组成部分。问题并不在于这些机制是否重要，而在于它们没有被提升为一个更底层的研究对象。

要解决这一点，课题目标必须从“做一个生命科学助手”转向“研究高代价科学工作流如何被控制”。一旦采用新的表述，系统的核心问题就不再是“如何理解生命科学任务并给出结果”，而变成了“在一个代价高、链路长、存在失败与恢复、且允许人工介入的科学工作流中，系统应依据什么原则决定下一步如何推进”。这样一来，系统关注的重点就会发生变化：难度、风险、证据充分性、恢复空间和人工介入时机，不再只是辅助性的描述概念，而会成为驱动执行策略变化的核心变量。

这种重新定义同时解决了三个问题。

首先，它能够把课题从一个应用层面的生命科学助手，提升为一个更通用的计算机问题。只要研究对象被定义为“高代价科学工作流控制”，那么蛋白设计就不再是课题本体，而只是一个验证场景。这样做的意义在于，系统的核心方法不仅可以服务于蛋白设计，也可能扩展到其他高成本科研流程，例如复杂文献调研、材料设计或多阶段实验分析。课题不再被某个具体应用域锁死，而开始具备“方法先于场景”的结构。

其次，这种重新定义有助于切断与“生命科学助手”父集之间的包含关系。如果别人的系统关注的是“如何面向生命科学任务提供智能帮助”，而当前课题关注的是“如何控制高代价工作流的执行、恢复与人工介入”，那么两者的层级已经不同。前者关注任务入口和任务处理，后者关注执行治理逻辑。此时，生命科学助手反而可以被看作该控制方法的上层入口，而不是其父集。

再次，从计算机专业毕设的角度看，这种重新定义更容易形成明确的算法主线。因为一旦将问题定义为“工作流控制”，就必须回答一系列典型的计算机问题：系统状态如何表示，哪些变量是不可直接观测的，运行时观测如何更新系统判断，动作空间如何定义，目标函数如何刻画成本、风险和收益之间的平衡，以及在有限预算和不确定性条件下系统如何做出决策。这些问题天然带有控制、调度、约束优化和序贯决策的特征，明显比“继续扩展生命科学助手能力”更接近一个可论证的研究方向。

基于这一判断，新的选题目标可以概括为：\textbf{不再把毕设定义为生命科学多智能体助手，而是定义为面向高代价科学工作流的运行时治理与控制问题，并以蛋白设计场景作为验证对象。} 在这一目标下，现有系统的角色也会发生变化。它不再是毕设本身的全部内容，而是一个实验平台和运行载体。状态机、人工介入、恢复链、候选生成和实验框架，都不再只是零散功能，而成为后续控制算法可依附的基础环境。

从时间角度看，这种重新定义同样具有现实意义。由于剩余周期已经不足两个月，如果继续沿“做更完整的生命科学助手”这一方向推进，项目边界只会不断扩大，论文主线却难以收敛。相反，如果此时将目标收缩到“高代价科学工作流的控制”这一更聚焦的问题上，后续工作就更容易围绕一个明确主线展开：系统状态如何抽象，难度和风险如何进入决策闭环，控制器如何决定继续执行、恢复、人工介入或终止，以及这种控制相对静态流程究竟能带来什么收益。

\paragraph{小结}
因此，当前毕设需要完成一次本质上的选题重构，即从“生命科学多智能体助手系统”转向“面向高代价科学工作流的运行时控制问题”。蛋白设计仍是最重要的应用验证场景，但它不再是课题定义本身。只有完成这一层转换，后续关于控制器、信念状态建模、创新点与可行性的讨论才真正成立。
